\documentclass[11pt]{article}

\usepackage[margin=1in]{geometry}
\usepackage{amsmath,amssymb,amsthm,mathtools}
\usepackage{booktabs,longtable,array}
\usepackage[T1]{fontenc}
\usepackage{lmodern}
\usepackage{microtype}
\usepackage{xcolor}
\usepackage{hyperref}
\usepackage{xurl}

\hypersetup{
  colorlinks=true,
  linkcolor=blue!55!black,
  citecolor=blue!55!black,
  urlcolor=blue!55!black,
  pdftitle={Audit of a rank-four counterexample to Grothendieck's killed-by-order question}
}

\newtheorem{theorem}{Theorem}[section]
\newtheorem{proposition}[theorem]{Proposition}
\newtheorem{lemma}[theorem]{Lemma}
\newtheorem{corollary}[theorem]{Corollary}
\theoremstyle{remark}
\newtheorem{remark}[theorem]{Remark}

\newcommand{\F}{\mathbf F}
\newcommand{\Z}{\mathbf Z}
\newcommand{\Spec}{\operatorname{Spec}}
\newcommand{\gr}{\operatorname{gr}}
\newcommand{\ord}{\operatorname{ord}}
\newcommand{\id}{\operatorname{id}}
\newcommand{\length}{\operatorname{length}}
\newcommand{\q}{\mathfrak q}
\newcommand{\m}{\mathfrak m}
\newcommand{\eps}{\varepsilon}

\title{Audit of a Rank-Four Counterexample to\\
Grothendieck's Killed-by-Order Question}
\author{Independent reconstruction and verification of the workspace construction}
\date{10 July 2026}

\begin{document}
\maketitle

\begin{abstract}
This document audits the claim that a previous computational pass found a
finite locally free group scheme of rank four which is not killed by four.
The claim is confirmed.  We give the group scheme as a completely specified
free rank-four Hopf algebra over a finite Artin local ring, prove that the
bialgebra equations and automatic-antipode argument are sufficient, and
explain the exact filtered linear-algebra certificates proving that two
coordinates of the fourth-power word are nonzero.  The base has
characteristic exactly \(4\), maximal ideal of fourth power zero, embedding
dimension \(15\), Hilbert function \((1,15,107,509)\), and length \(632\).
The group law is noncommutative.  Its fourth-power defect is a nonzero,
third-order socle-valued map, while its eighth-power word is trivial.

The result is a computer-assisted disproof of Grothendieck's general
killed-by-order assertion.  The decisive calculations are exact finite
linear algebra over \(\F_2\): no timeout, probabilistic computation, or
heuristic Groebner remainder is used as evidence.
\end{abstract}

\tableofcontents

\section{Verdict and scope}

Grothendieck asked whether a finite locally free group scheme of constant
rank \(n\) over an arbitrary base is killed by its order: equivalently,
whether the pointwise \(n\)-th power word factors through the unit section.
The question is recorded in SGA~3, Expose VIII, Remarque 7.3.1; convenient
modern accounts are Schoof's survey and Torti's recent paper
\cite{SGA3,SchoofSurvey,Torti}.  For noncommutative group schemes this power
word need not be a group homomorphism, but it is still a morphism of schemes.

The audited construction gives a single explicit pair
\[
  G=\Spec A \longrightarrow \Spec B
\]
with \(A\) finite free of rank \(4\) over \(B\), such that
\[
  [4]^\#\ne\eta\eps:A\longrightarrow A.
\]
Thus the statement is negatively resolved: the general conjectural assertion
is false.  This conclusion is stronger than finding a formal obstruction or
a non-liftable infinitesimal seed; the universal cubic quotient used below is
itself an honest finite base ring and carries the honest group scheme.

The construction is explicit but not small.  It is best viewed as a
universal third-order deformation quotient.  Nothing in the audit proves
that this base, its length, or its embedding dimension is minimal.

\section{The power word in Hopf-algebra coordinates}

Let \(R\) be a commutative ring and let \(G=\Spec A\) be a finite locally
free affine group scheme over \(R\).  The coordinate algebra \(A\) is a
commutative \(R\)-algebra even when the group law on \(G\) is noncommutative;
noncommutativity is encoded by failure of cocommutativity of the coproduct
\(\Delta:A\to A\otimes_R A\).

Write \(m:A\otimes_R A\to A\) for multiplication and put
\[
  \phi=m\circ\Delta=[2]^\#.
\]
The square of an element can be squared again without assuming that the group
is commutative, so
\[
  [4]^\#=\phi\circ\phi.
\]
If \(I_A=\ker(\eps:A\to R)\), then \(G\) is killed by four precisely when
\[
  [4]^\#(I_A)=0,
\]
or equivalently when \([4]^\#=\eta\eps\).

\section{The explicit universal cubic group scheme}

\subsection{The coefficient ring and its variables}

Let
\[
 P=\Z_{(2)}\bigl[m_{ij}^{\,r},c_{ijk}\bigr],
\]
where
\[
  1\le i\le j\le3,\quad 1\le r\le3,
  \qquad 1\le i,j,k\le3.
\]
There are \(6\cdot3=18\) multiplication variables and \(3^3=27\)
reduced-coproduct variables, hence \(45\) variables in total.  Set
\[
  \q=(2,m_{ij}^{\,r},c_{ijk})\subset P.
\]

On the free \(P\)-module
\[
  A_P=P\{1,e_1,e_2,e_3\}
\]
define a commutative, unital, augmented multiplication by
\begin{equation}\label{eq:multiplication}
 e_i e_j=
 \begin{cases}
   e_3,&\{i,j\}=\{1,2\},\\
   0,&\text{otherwise}
 \end{cases}
 +\sum_{r=1}^3m_{ij}^{\,r}e_r,
 \qquad 1\le i\le j\le3,
\end{equation}
and \(\eps(e_i)=0\).  Define a counital coproduct by
\begin{equation}\label{eq:coproduct}
\begin{split}
 \Delta(e_i)={}&e_i\otimes1+1\otimes e_i\\
 &+\delta_{i3}(e_1\otimes e_2+e_2\otimes e_1)
 +\sum_{j,k=1}^3c_{ijk}e_j\otimes e_k,
\end{split}
\end{equation}
and \(\Delta(1)=1\otimes1\).

\subsection{The 189 bialgebra equations}

For completeness, the equations can be specified without printing their
large expanded polynomial forms.  Put \(e_0=1\), and define total structure
constants \(M_{ab}^{\,r}\) and \(D_i^{\,ab}\) by
\[
 e_a e_b=\sum_{r=0}^3M_{ab}^{\,r}e_r,
 \qquad
 \Delta(e_i)=\sum_{a,b=0}^3D_i^{\,ab}e_a\otimes e_b,
\]
using \eqref{eq:multiplication} and \eqref{eq:coproduct}.  In particular,
\(M_{0b}^{\,r}=\delta_{br}\), \(M_{a0}^{\,r}=\delta_{ar}\), and
\(M_{ab}^{\,0}=0\) for \(a,b>0\).

Let \(I\subset P\) be generated by the following coefficients:
\begin{align}
 A_{ijk}^{\,r}
 &=\sum_{s=0}^3M_{ij}^{\,s}M_{sk}^{\,r}
   -\sum_{s=0}^3M_{jk}^{\,s}M_{is}^{\,r},
 &&1\le i,j,k,r\le3,\label{eq:assoc}\\
 C_{ij}^{\,rs}
 &=\sum_{t=0}^3M_{ij}^{\,t}D_t^{\,rs}
   -\sum_{a,b,c,d=0}^3D_i^{\,ab}D_j^{\,cd}
       M_{ac}^{\,r}M_{bd}^{\,s},
 &&1\le i\le j\le3,\quad1\le r,s\le3,\label{eq:compat}\\
 K_i^{\,abc}
 &=\sum_{r=0}^3D_i^{\,rc}D_r^{\,ab}
   -\sum_{s=0}^3D_i^{\,as}D_s^{\,bc},
 &&1\le i,a,b,c\le3.\label{eq:coassoc}
\end{align}

After zero polynomials and exact duplicates are removed, the three families
contribute respectively \(54\), \(54\), and \(81\) polynomials, for a total
of \(189\).

The equations with a \(1=e_0\) leg have not been omitted as assumptions.
They hold identically from the augmented multiplication and counital normal
form.  Thus \eqref{eq:assoc}--\eqref{eq:coassoc} impose the complete
associativity, coproduct-multiplicativity, and coassociativity axioms.

Now define
\begin{equation}\label{eq:B-definition}
  B=P/(I+\q^4),\qquad A=A_P\otimes_P B,
  \qquad G=\Spec A.
\end{equation}
This is the promised explicit object.

\subsection{Closed fiber}

At the closed point \(\q\), all \(m\)'s and \(c\)'s vanish and \(2=0\).  The
special fiber has
\[
  e_1^2=e_2^2=e_3^2=e_1e_3=e_2e_3=0,
  \qquad e_1e_2=e_3,
\]
and
\[
\begin{aligned}
 \Delta(e_1)&=e_1\otimes1+1\otimes e_1,\\
 \Delta(e_2)&=e_2\otimes1+1\otimes e_2,\\
 \Delta(e_3)&=e_3\otimes1+1\otimes e_3
              +e_1\otimes e_2+e_2\otimes e_1.
\end{aligned}
\]
Consequently
\[
  A\otimes_B\F_2\simeq\F_2[x,y]/(x^2,y^2),
  \qquad(e_1,e_2,e_3)=(x,y,xy),
\]
with the usual Hopf structure of \(\alpha_2\times\alpha_2\).

\section{Why this is an actual finite flat group scheme}

\subsection{The base is finite Artin local}

Every generator of \(I\) belongs to \(\q\), because its reduction at the
closed point is a valid bialgebra equation over \(\F_2\).  Hence
\[
  \q^4\subset I+\q^4\subset\q,
  \qquad \sqrt{I+\q^4}=\q.
\]
It follows that \(B\) is local, with residue field \(\F_2\), and its maximal
ideal
\[
  \m=\q B
\]
satisfies \(\m^4=0\).  The quotient \(P/\q^4\) is finite, so \(B\) is a
finite Artin local ring.

The exact filtered ranks computed below imply
\begin{equation}\label{eq:hilbert}
 \dim_{\F_2}\gr_{\m}(B)=(1,15,107,509).
\end{equation}
Therefore
\[
 \length(B)=632,\qquad |B|=2^{632},
\]
and the embedding dimension is \(15\).  Since the degree-three piece is
nonzero, the Loewy length is exactly \(4\).  Moreover \(\m^3\) is a
\(509\)-dimensional subspace of the socle, so this universal base is very far
from Gorenstein.

An additional exact membership audit gives
\begin{equation}\label{eq:char4}
  2\ne0\quad\text{and}\quad4=0\quad\text{in }B.
\end{equation}
Thus \(B\) has characteristic exactly \(4\).  The certificate for
\eqref{eq:char4} is described in Section~\ref{sec:extra-certificates}.

\subsection{Rank four is built into the construction}

The equations in \(I\) constrain structure constants; they do not quotient
the underlying module \(B\{1,e_1,e_2,e_3\}\).  Hence \(A\) is free of rank
four as a \(B\)-module.  Equations \eqref{eq:assoc}--\eqref{eq:coassoc}
make it a finite free \(B\)-bialgebra.

\subsection{The antipode is automatic}

Let
\[
  \mathcal C=\operatorname{End}_B(A)
\]
with convolution product and unit \(\eta\eps\).  This is a finite \(B\)-algebra.
Modulo \(\m\), the element \(\id_A\) is a convolution unit because the closed
fiber \(\alpha_2^2\) is Hopf.  The ideal \(\m\mathcal C\) is nilpotent, so a
unit modulo \(\m\mathcal C\) lifts to a unit in \(\mathcal C\).  Explicitly,
the error from a lifted inverse lies in \(\m\mathcal C\) and is inverted by a
finite geometric series.

The convolution inverse of \(\id_A\) is an antipode.  Therefore \(A\) is a
Hopf algebra and
\[
  G\longrightarrow\Spec B
\]
is a finite locally free group scheme of rank four.  No unverified antipode
equations are missing from the construction.

\section{The exact filtered certificate}

\subsection{The mixed \texorpdfstring{\(\q\)}{q}-adic grading}

Index the \(45\) deformation variables as \(p_0,\ldots,p_{44}\), as in
Appendix~\ref{app:numbering}.  For a nonzero integer term \(c p^\alpha\), put
\[
  \ord_{\q}(c p^\alpha)=v_2(c)+|\alpha|.
\]
Then
\[
 \gr_{\q}P\simeq\F_2[\tau,p_0,\ldots,p_{44}],
 \qquad \tau=\operatorname{in}_{\q}(2).
\]
The variable \(\tau\) records genuine \(2\)-adic carries; it is not an
independent parameter of the exact ring.  Integer coefficients are combined
before \(v_2\) is taken.  Thus two equal odd terms produce a term containing
\(\tau\), rather than cancelling erroneously in advance.

\subsection{The filtered-generation lemma}

The previous note stated the recursion with a hypothesis that was too weak.
The following is the precise invariant actually maintained by both
implementations.

\begin{lemma}[Exact filtered recursion]\label{lem:filtered}
Suppose a finite list \(C_d\subset I\cap\q^d\) generates
\(I\cap\q^d\) as a \(P\)-module.  Row-reduce the degree-\(d\) initial forms
over \(\F_2\), retaining exact polynomial combinations.  Let \(G_d\) be exact
lifts of a basis of the initial row space, and let \(H_d\) be exact residuals
of a basis of all constant \(\F_2\)-linear relations.  Then
\[
 C_{d+1}=
 \{2g,p_0g,\ldots,p_{44}g:g\in G_d\}\cup H_d
\]
generates \(I\cap\q^{d+1}\) as a \(P\)-module.
\end{lemma}

\begin{proof}
The exact row operations are triangular elementary additions.  Consequently
\(G_d\cup H_d\) generates the same \(P\)-module as \(C_d\).  The degree-\(d\)
initial forms of the elements of \(G_d\) are linearly independent, while
\(H_d\subset\q^{d+1}\).

Take \(F\in I\cap\q^{d+1}\) and write
\[
 F=\sum_g a_g g+\sum_h b_h h,
 \qquad g\in G_d,\quad h\in H_d.
\]
Independence of the degree-\(d\) initials forces every \(a_g\) to lie in
\(\q\).  Since \(\q=(2,p_0,\ldots,p_{44})\), each \(a_g g\) is generated by
the displayed \(46\) multiples of \(g\), while the \(h\)'s are retained.
The reverse containment is immediate.
\end{proof}

At \(d=1\), the \(189\) original equations generate \(I=I\cap\q\), so the
lemma applies inductively.  It also remains valid after localization from
\(\Z[p]\) to \(\Z_{(2)}[p]\): every odd denominator is a unit with initial
form \(1\), and its correction from \(1\) lies in \(\q\).  This supplies the
localization argument omitted from the earlier note.

Terms of \(\q\)-order at least \(4\) can safely be discarded.  Addition and
multiplication cannot lower \(\q\)-order, so such terms can never return to a
degree relevant modulo \(\q^4\).

\subsection{The three exact stages}

The independently reproduced ranks are:
\begin{center}
\begin{tabular}{@{}rrrrr@{}}
\toprule
degree & ambient dimension & candidates & nonzero candidates & rank\\
\midrule
1 & \(46\)    & \(189\)   & ---      & \(31\)\\
2 & \(1081\)  & \(1584\)  & ---      & \(974\)\\
3 & \(17296\) & \(45414\) & \(45402\)& \(16787\)\\
\bottomrule
\end{tabular}
\end{center}

At degree one the kernel has dimension \(189-31=158\).  Therefore the
degree-two list has
\[
  46\cdot31+158=1584
\]
elements.  Its kernel has dimension \(1584-974=610\), so Lemma
\ref{lem:filtered} gives exactly
\[
  46\cdot974+610=45414
\]
cubic candidates.  This count is structural: the exhaustive recursion
produces all \(45{,}414\) candidates (of which \(45{,}402\) have nonzero cubic
initial), and their span is \(\gr_3(I)\).

The Hilbert function \eqref{eq:hilbert} follows at once:
\[
 1,\quad46-31=15,\quad1081-974=107,
 \quad17296-16787=509.
\]

\section{The fourth-power calculation}

Write
\[
  \phi(e_i)=\sum_{r=1}^3\phi_i^{\,r}e_r.
\]
There is no scalar coordinate because \(\phi\) preserves augmentation.  The
nine target polynomials are
\begin{equation}\label{eq:targets}
  T_{ir}=\sum_{s=1}^3\phi_i^{\,s}\phi_s^{\,r},
  \qquad1\le i,r\le3,
\end{equation}
which are precisely the \(e_r\)-coordinates of
\([4]^\#(e_i)=\phi^2(e_i)\).  The source orders them as
\[
  3(i-1)+(r-1),
\]
so targets \(2\) and \(5\) are \(T_{13}\) and \(T_{23}\).

All nine targets have \(\q\)-order at least two.  After exact quadratic
cancellation, the cubic normal remainders of targets \(2\) and \(5\) are
\begin{align}
 \rho_2
 &=p_0p_{18}p_{19}+p_0p_{18}p_{21}
   =m_{11}^{\,1}c_{111}(c_{112}+c_{121}),\label{eq:rho2}\\
 \rho_5
 &=p_0p_{18}p_{28}+p_0p_{18}p_{30}
   =m_{11}^{\,1}c_{111}(c_{212}+c_{221}).\label{eq:rho5}
\end{align}
The term \emph{normal} here refers to the specified high-pivot ordering; the
nonzero quotient classes are intrinsic, while these representatives depend
on the reduction convention.

Two explicit \(\F_2\)-linear functionals \(L_2,L_5\) on the \(17296\)-dimensional
cubic space have the following properties:
\[
\begin{array}{c|cc}
 &\rho_2&\rho_5\\\hline
 L_2&1&0\\
 L_5&0&1
\end{array}
\]
and both functionals annihilate every one of the \(45414\) exhaustive cubic
ideal candidates.  Each has support \(38\); the full supports are printed in
Appendix~\ref{app:duals}.  For the hashes, monomials are ordered
lexicographically as nondecreasing triples, the coefficient bits are packed
as a minimal-length little-endian byte string, and SHA-256 is applied to that
byte string.  The resulting hashes are
\begin{align*}
L_2:\;&\texttt{1d8738d672e98ddd690657b14821e52e2}\\
     &\texttt{d2445551a9e9d5b556440d66eac527f},\\[2pt]
L_5:\;&\texttt{505bbba3120840f5644bf5b0b4ec0fde}\\
     &\texttt{df0861cd702fdddb6240fdad84f11dbd}.
\end{align*}
Therefore neither residual belongs to \(\gr_3(I)\), and hence
\[
  T_{13},T_{23}\notin I+\q^4.
\]
The other seven target rows reduce to zero through degree three; exact
filtered lifting therefore puts them in \(I+\q^4\).

\begin{theorem}[Audited counterexample]\label{thm:main}
For \(B,A,G\) defined in \eqref{eq:B-definition}, the morphism
\(G\to\Spec B\) is a finite locally free group scheme of rank \(4\), and
\[
  [4]^\#\ne\eta\eps.
\]
More precisely,
\begin{equation}\label{eq:p4-description}
\begin{aligned}
 [4]^\#(e_1)&=\theta_1e_3,\\
 [4]^\#(e_2)&=\theta_2e_3,\\
 [4]^\#(e_3)&=0,
\end{aligned}
\end{equation}
where \(\theta_1,\theta_2\in\m^3\) are nonzero and linearly independent in
\(\gr_3(B)=\m^3\).  Their normal representatives are
\eqref{eq:rho2} and \eqref{eq:rho5}.
\end{theorem}

\begin{proof}
The finite-flat and Hopf assertions were proved in Section~4.  The seven zero
target coordinates and the two nonzero target coordinates give
\eqref{eq:p4-description}.  Since \(A\) is a free \(B\)-module with basis
\(1,e_1,e_2,e_3\), a nonzero coefficient gives a nonzero \(B\)-linear map.
The two dual functionals give the stated independence.
\end{proof}

\section{A concrete description of the resulting group scheme}

The example has the following useful features.
\begin{enumerate}
\item Its unique geometric fiber is \(\alpha_2^2\), and is killed by two.
The failure of killedness is therefore purely infinitesimal.

\item The fourth-power defect lands in the line \(Be_3\subset I_A\) and is
scaled by two independent elements of \(\m^3\).  Since \(\m\m^3=0\), the
defect is socle-valued and invisible through the first two infinitesimal
layers.

\item The total coordinate algebra is the free module with multiplication
table \eqref{eq:multiplication}; it is not simply
\(B[x,y]/(x^2,y^2)\).  The latter describes only the closed fiber.

\item The group law is genuinely noncommutative.  Indeed, the coefficient of
\(e_1\otimes e_2-e_2\otimes e_1\) in \(\Delta(e_1)\) is
\[
  c_{112}-c_{121},
\]
whose nonzero linear initial class survives the \(31\)-dimensional linear
equation space.  An explicit two-term dual certificate is recorded in
Section~\ref{sec:extra-certificates}.  Thus \(\Delta\) is not cocommutative.
This is necessary: Deligne's theorem rules out a counterexample among
commutative group schemes \cite{SchoofSurvey}.

\item The group scheme is killed by eight.  Reduction modulo \(\m\) gives
\[
  [2]^\#(I_A)\subset\m I_A.
\]
By Theorem~\ref{thm:main},
\([4]^\#(I_A)\subset\m^3I_A\).  Therefore
\[
 [8]^\#(I_A)
 =[2]^\#\bigl([4]^\#(I_A)\bigr)
 \subset\m^4I_A=0.
\]
In the pointwise-word sense, its exact power exponent is \(8\): it is killed
by eight but not by four.
\end{enumerate}

\section{Supplementary exact certificates}\label{sec:extra-certificates}

\subsection{Characteristic four}

Exact cancellation of the degree-one initial of \(2\) by the \(31\) linear
pivot lifts leaves, through relevant order,
\begin{align*}
R_2={}&-3p_0p_{18}+p_1p_{27}+p_2p_{36}
 -2p_3p_{19}-2p_3p_{21}\\
&-2p_6p_{20}-2p_6p_{24}.
\end{align*}
Its degree-two normal remainder is
\[
  p_0p_{18}+p_1p_{27}.
\]
The functional supported on
\[
 \{(1,27),(10,31),(10,41),(10,43),
   (14,31),(14,41),(14,43)\}
\]
annihilates all \(1584\) exhaustive quadratic ideal candidates and evaluates
to \(1\) on this residual.  Its bit-vector hash is
\[
\begin{split}
&\texttt{19adb92126a6c8f844cb7d4b217e7a099}\\
&\texttt{dadbd4cfe0025cea2ea0b4419e60478}.
\end{split}
\]
Thus \(2\notin I+\q^4\).  In contrast, exact degree-two cancellation of
\(4\) leaves
\[
 -6p_0p_{18}+2p_1p_{27}+2p_2p_{36},
\]
whose degree-three initial cancels against the exhaustive cubic list, leaving
only \(\q^4\).  Hence \(4\in I+\q^4\), proving \eqref{eq:char4}.

\subsection{Noncocommutativity}

The linear initial of \(c_{112}-c_{121}=p_{19}-p_{21}\) is
\(p_{19}+p_{21}\).  It is unchanged by reduction modulo the degree-one
initial ideal.  The linear functional supported on \(p_{21}\) and \(p_{41}\)
annihilates all \(189\) equation initials and evaluates to \(1\) on
\(p_{19}+p_{21}\).  Therefore
\[
  c_{112}-c_{121}\notin I+\q^4,
\]
so the coproduct in the constructed Hopf algebra is not cocommutative.

\section{Independent reproduction and audit trail}

\subsection{Source and certificate paths}

The principal files are:
\begin{center}
\begin{tabular}{@{}p{0.42\textwidth}p{0.49\textwidth}@{}}
\toprule
path & role\\
\midrule
\path{m2/universal_local_rank4.m2}
  & constructs and exports the \(189\) equations and nine targets\\
\path{scripts/independent_audit_mixed_a2a2_export.py}
  & independently parses the exact M2 export and performs filtered elimination\\
\path{scripts/audit_universal_rank4_quadratic.py}
  & independently reconstructs the chart and supplies the first certificate path\\
\path{scripts/compare_universal_rank4_m2_export_20260710.py}
  & compares every relevant exact M2 and Python coefficient\\
\path{m2/verify_mixed_a2a2_dual_direct.m2}
  & native M2 matrix verification of the explicit dual supports\\
\path{scripts/audit_rank4_counterexample_base_invariants_20260710.py}
  & verifies characteristic four, the Hilbert function, and noncocommutativity\\
\path{logs/rank4_counterexample_audit_20260710.txt}
  & persistent command/output transcript\\
\bottomrule
\end{tabular}
\end{center}

\subsection{Fresh commands and results}

The export was regenerated with Macaulay2 1.25.11:
\begin{verbatim}
env RANK4_BRANCH=0 RANK4_INTEGRAL=1 RANK4_EXPORT_POLYS=1 \
    M2 --script m2/universal_local_rank4.m2 \
    > /tmp/m2_rank4_export_audit_20260710.txt
\end{verbatim}
It contained \(189\) equation records and nine target records and had the
same SHA-256 as the frozen prior export:
\begin{verbatim}
686c35c91b74af2a414dd119024411266fb0056865dd1d19676c20485f93b6b7
\end{verbatim}
The M2/Python coefficient comparison reported exact equality through all
terms relevant modulo \(\q^4\), with combined canonical hash
\begin{verbatim}
2f41a12591ea296a33cd7589dd4f8911ea108dd04e957ba6f516a26a90aca695
\end{verbatim}
The independent Python elimination reproduced the rank table, both
nonmembers, both \(38\)-term supports, and both hashes in under one second,
using about \(63\) MB peak RSS.

The native Macaulay2 verifier was a newly added path which had not previously
been logged as executed.  During this audit it exposed two syntax defects
before reaching any mathematical calculation: the auxiliary monoid needed
to be constructed before extending scalars, and two list-concatenation
operators could not begin a new parsed statement.  The syntax-only repairs
did not alter any equation, row, target, or dual support.  The repaired
verifier then completed with exit code zero:
\begin{verbatim}
degree1 ... rank=31 relations=158
degree2 ... rank=974 relations=610
degree3 ... candidates=45414 L2_violations=0 L5_violations=0
target2_dual_value=1 target5_dual_value=1
AUDIT PASS: exact M2 filtered certificate through q^4
\end{verbatim}
It took \(53.8\) seconds and peaked at approximately \(964\) MiB RSS.  Thus
the native certificate run remained comfortably bounded on the audit Mac.

\subsection{Frozen and audited hashes}

\begingroup\scriptsize
\begin{longtable}{@{}p{0.69\textwidth}p{0.25\textwidth}@{}}
\toprule
SHA-256 & file\\
\midrule
\endhead
\texttt{ad55eab4fa40aa31cd39fcccbbfc7f6f97fd2f2aa01398ffc8c3cc7cf3546b89}
& \path{m2/universal_local_rank4.m2}\\
\texttt{43d0be4f325a599a530bd87cc3653b119fae23a9ad61a2cb67b24ad975f32522}
& independent Python checker\\
\texttt{73c7befdb01a23c8622dc58d61fc77bc78c8566437b868e1507905b070f2207f}
& first Python reconstruction\\
\texttt{ede087dddb8ab5a510375f6ed71eae2e68eab89774666e079fa91a9954b5e7c0}
& M2/Python comparison script\\
\texttt{b278ec35fcb252d37b446370089233045fa7cbb6c5d8efd8fa2393f1dd7c5622}
& repaired native M2 verifier\\
\texttt{cc2a89eaa66407d27dd89c1287177d3a68d92503252b2f5700a217224535cb48}
& base-invariant checker\\
\bottomrule
\end{longtable}
\endgroup

\section{Audit qualifications}

The result passes the audit, subject to the following precise qualifications.
\begin{enumerate}
\item This is an exact computer-assisted proof.  The \(38\)-term
functionals and the filtered-generation lemma make the finite computation
auditable, but the \(45{,}414\) row evaluations are delegated to the supplied
programs.

\item The earlier written filtered lemma was false under the weak assumption
that a candidate list merely spans the degree-\(d\) initials.  Lemma
\ref{lem:filtered} states the stronger module-generation invariant actually
maintained by the code and repairs the exposition.

\item The displayed cubic normal forms are canonical only relative to the
specified variable order and high-pivot reduction.  Their nonzero classes
and the failure of killedness are independent of that convention.

\item The workspace is not a Git worktree.  Provenance therefore rests on
the hashes above and on byte-for-byte regeneration of the M2 export.

\item Several older handoff bodies say that no counterexample was found.
Their opening supersession banners point to the later cubic audit.  Earlier
rank-\(16\) and pinned rank-\(4\) apparent seeds were explicitly retracted for
unrelated scanner or pinning errors; they are not evidence for the present
construction.

\item The conclusion is a mathematical audit of a local artifact, not a
claim about publication priority or completion of external peer review.
As late as Torti's article published in December 2025 (and its March 2026
arXiv revision), the general question was described as open and Schoof's
theorem as the best general positive result
\cite{Schoof2001,Torti}.
\end{enumerate}

\section{Conclusion}

The previous agent's decisive claim is correct.  The ring
\[
 B=\Z_{(2)}[m_{ij}^{\,r},c_{ijk}]/(I+\q^4)
\]
and the multiplication and coproduct tables
\eqref{eq:multiplication}--\eqref{eq:coproduct} define a finite locally free
rank-four group scheme \(G/\Spec B\).  Exact dual certificates show that the
\(e_3\)-coordinates of \([4]^\#(e_1)\) and \([4]^\#(e_2)\) are nonzero.
Therefore \(G\) is not killed by four.  It is a characteristic-four,
noncommutative, purely infinitesimal counterexample whose fourth-power defect
appears only in the third-order socle; it is nevertheless killed by eight.

Consequently Grothendieck's killed-by-order assertion over arbitrary bases is
false.

\appendix

\section{Variable numbering}\label{app:numbering}

For the unordered pairs
\[
 (1,1),(1,2),(1,3),(2,2),(2,3),(3,3)
\]
in that order, let \(\rho(i,j)\in\{0,1,2,3,4,5\}\) denote the position.
Then
\[
 p_{3\rho(i,j)+r-1}=m_{ij}^{\,r}
 \quad(1\le r\le3),
\]
and
\[
 p_{18+9(i-1)+3(j-1)+(k-1)}=c_{ijk}.
\]
In particular,
\[
\begin{gathered}
p_0=m_{11}^{\,1},\quad p_{18}=c_{111},\quad
p_{19}=c_{112},\quad p_{21}=c_{121},\\
p_{28}=c_{212},\quad p_{30}=c_{221}.
\end{gathered}
\]
The forty-sixth associated-graded variable is \(\tau=\operatorname{in}(2)\);
the explicit cubic duals below happen to have no \(\tau\)-supported term.

\section{The two explicit cubic dual supports}\label{app:duals}

For a set \(S\) of nondecreasing triples in \(\{0,\ldots,45\}\), let
\(L_S\) be the \(\F_2\)-linear functional which sums the coefficients of the
corresponding cubic monomials.  The target-\(2\) support is
\begingroup\small
\begin{align*}
S_2=\{&
(0,18,21),(0,18,41),(0,21,30),(0,21,38),(0,30,31),\\
&(0,30,43),(0,31,42),(0,38,41),(0,41,42),(0,42,43),\\
&(1,18,30),(1,18,42),(1,30,30),(1,30,38),(1,38,42),\\
&(1,42,42),(8,18,21),(8,18,41),(8,21,30),(8,21,38),\\
&(8,30,31),(8,30,43),(8,31,42),(8,38,41),(8,41,42),\\
&(8,42,43),(10,21,21),(10,21,31),(10,21,43),(10,31,41),\\
&(10,41,41),(10,41,43),(14,21,21),(14,21,31),(14,21,43),\\
&(14,31,41),(14,41,41),(14,41,43)\}.
\end{align*}
\endgroup

The target-\(5\) support is
\begingroup\small
\begin{align*}
S_5=\{&
(0,18,30),(0,18,42),(0,30,30),(0,30,38),(0,38,42),\\
&(0,42,42),(8,18,30),(8,18,42),(8,30,30),(8,30,38),\\
&(8,38,42),(8,42,42),(9,21,21),(9,21,31),(9,21,43),\\
&(9,31,41),(9,41,41),(9,41,43),(10,18,21),(10,18,41),\\
&(10,21,30),(10,21,38),(10,30,31),(10,30,43),(10,31,42),\\
&(10,38,41),(10,41,42),(10,42,43),(14,18,21),(14,18,41),\\
&(14,21,30),(14,21,38),(14,30,31),(14,30,43),(14,31,42),\\
&(14,38,41),(14,41,42),(14,42,43)\}.
\end{align*}
\endgroup

The native M2 verifier evaluates \(L_{S_2}\) and \(L_{S_5}\) on each of the
\(45414\) cubic candidates.  Both violation counts are zero, and their values
on the respective exactly lifted target residuals are one.

\begin{thebibliography}{9}

\bibitem{SGA3}
M.~Demazure and A.~Grothendieck (eds.),
\emph{Schemas en groupes (SGA 3)},
Exposes VIIA and VIII, Lecture Notes in Mathematics, Springer.

\bibitem{Schoof2001}
R.~Schoof,
\emph{Finite flat group schemes over local Artin rings},
Compositio Mathematica \textbf{128} (2001), 1--15.
\url{https://doi.org/10.1023/A:1017560215203}.

\bibitem{SchoofSurvey}
R.~Schoof,
\emph{Is a finite locally free group scheme killed by its order?},
survey/preliminary version, 2017.
\url{https://www.mat.uniroma2.it/~schoof/schoof_oortAAG.pdf}.

\bibitem{Torti}
E.~Torti,
\emph{Lagrange's theorem for a class of finite flat group schemes over local
Artin rings},
Documenta Mathematica (2025), published online first,
\url{https://doi.org/10.4171/DM/1055}; see also the revised arXiv:2411.12129.
\url{https://arxiv.org/abs/2411.12129}.

\end{thebibliography}

\end{document}
